% ============================================================================
%  ESTADÍSTICA DESCRIPTIVA — Sesión 1
%  Semana 1: 10 al 15 de agosto de 2026
%  Presentación del curso + 1.1 Conceptos básicos: población, muestra,
%  variable y dato
% ============================================================================
\documentclass[aspectratio=169,10pt]{beamer}

\usepackage[utf8]{inputenc}
\IfFileExists{spanish.ldf}{\usepackage[spanish,es-noshorthands]{babel}}{}
\usepackage{amsmath,amssymb}
\usepackage{booktabs}
\usepackage{tikz}
\usepackage{listings}

\usetheme{Uniminuto}

\uniasignatura{Estadística descriptiva}
\uniprofesor{Andrés Felipe Puerta Velez}
\unicorreo{andres.puerta-v@uniminuto.edu.co}
\uninrc{CBASBL031 - CBASBL141}
\unisemana{Sesión 1 — Semana 1 · 10 al 15 de agosto de 2026}


\newcommand{\fx}[1]{{\ttfamily\color{uniAzul}#1}}
\newcommand{\listatiny}{\setbeamerfont{itemize/enumerate body}{size=\tiny}\tiny}

% ------------------------------------------------------------- CÓDIGO R ----
\lstdefinestyle{Rmini}{
  language=R,
  basicstyle=\ttfamily\tiny,
  keywordstyle=\color{uniAzul}\bfseries,
  commentstyle=\color{uniGris}\itshape,
  stringstyle=\color{uniTitulo},
  showstringspaces=false, numbers=none, breaklines=true,
  columns=fullflexible,
  backgroundcolor=\color{uniCremaClaro},
  frame=leftline, framerule=1.2pt, rulecolor=\color{uniNaranja},
  xleftmargin=0.35cm, framexleftmargin=0.35cm,
  framexrightmargin=0.15cm, framextopmargin=2pt, framexbottommargin=2pt,
  aboveskip=0.12cm, belowskip=0.12cm
}
\lstdefinestyle{Rout}{
  basicstyle=\ttfamily\tiny\color{black},
  showstringspaces=false, numbers=none, breaklines=true,
  columns=fullflexible,
  backgroundcolor=\color{uniGris!10},
  frame=leftline, framerule=1.2pt, rulecolor=\color{uniGris},
  xleftmargin=0.35cm, framexleftmargin=0.35cm,
  framexrightmargin=0.15cm, framextopmargin=2pt, framexbottommargin=2pt,
  aboveskip=0.12cm, belowskip=0.12cm
}

\begin{document}

\portada

\begin{frame}{Orden del día}
  \begin{enumerate}
    \item Llamado a lista, presentación del grupo y elección del representante
    \item Para qué sirve la estadística: tres historias y tres preguntas
    \item Temario, cronograma, evaluación y acuerdos mínimos
    \item \azul{Tema 1.1: población, muestra, unidad de observación, variable y dato}
    \item Cómo se produce un dato y cómo se guarda: la matriz de datos
    \item De dónde salen los datos: censo, muestra y sesgos de recolección
    \item Cómo se lee mal una estadística; taller de diseño y cierre
  \end{enumerate}

  \vspace{0.25cm}
  {\footnotesize\color{uniGris} Sesión de 3 horas. Al final de hoy queda
  definido el \destacar{conjunto de datos} que acompañará todo el semestre.}
\end{frame}

\begin{frame}{¿Para qué sirve la estadística descriptiva?}
  Un conjunto de datos en bruto no dice nada. \destacar{Describir} es
  convertirlo en algo que se pueda leer, comparar y comunicar.

  \vspace{0.3cm}
  \begin{columns}[T,onlytextwidth]
    \begin{column}{0.47\textwidth}
      \begin{block}{Lo que haremos}
        \footnotesize
        \begin{itemize}
          \item Organizar datos en tablas
          \item Representarlos en gráficos
          \item Resumirlos en unos pocos números
          \item Interpretarlos y sustentar decisiones
        \end{itemize}
      \end{block}
    \end{column}
    \begin{column}{0.47\textwidth}
      \begin{exampleblock}{Hasta dónde llega}
        \footnotesize
        La estadística descriptiva \azul{resume lo que ya se observó}.\\[0.3em]
        No afirma nada sobre los casos que \azul{no} se midieron.\\[0.3em]
        Esa disciplina no es una limitación: es lo que hace que un informe sea
        \destacar{defendible}.
      \end{exampleblock}
    \end{column}
  \end{columns}
\end{frame}

\begin{frame}{Tres historias que hicieron la estadística}
  \scriptsize
  \begin{block}{1854 — John Snow y el cólera en Londres}
    \tiny
    Cuando estalla un brote en el Soho, la teoría dominante culpa al ``aire
    viciado''. Snow hace otra cosa: \azul{marca en un mapa dónde vivía cada
    muerto}. Las marcas se apiñan alrededor de una bomba de agua de la calle
    Broad. Convence a las autoridades de retirar la palanca de la bomba.
    \destacar{No tenía una teoría mejor: tenía los datos organizados.}
  \end{block}

  \vspace{0.06cm}
  \begin{block}{1858 — Florence Nightingale y el diagrama de la rosa}
    \tiny
    De vuelta de la guerra de Crimea, quiere convencer al gobierno británico
    de que la mayoría de los soldados no moría por heridas de combate sino por
    \azul{enfermedades evitables en los hospitales}. Sabe que nadie va a leer
    su tabla, así que la convierte en un gráfico circular por meses. El
    gráfico funciona: \destacar{cambió la política sanitaria del ejército}.
  \end{block}

  \vspace{0.06cm}
  \begin{block}{1890 — el censo de Estados Unidos}
    \tiny
    El censo anterior había tardado casi ocho años en tabularse. Hermann
    Hollerith diseña una máquina de tarjetas perforadas para contarlo, y su
    empresa terminaría siendo parte de lo que hoy es \azul{IBM}.
    \destacar{La informática nació, literalmente, de un problema
    estadístico.}
  \end{block}

  \vspace{0.06cm}
  {\tiny\color{uniGris} Las tres tienen la misma moraleja:
  \destacar{los datos no convencen solos}. Convencen cuando alguien los
  organiza, los representa y los interpreta. Eso es exactamente lo que vamos a
  aprender.}
\end{frame}

\begin{frame}{Cuatro preguntas que llegan a su escritorio}
  \tiny
  \begin{center}
  \begin{tabular}{@{}p{1.4cm}p{4.8cm}p{5.8cm}@{}}
    \toprule
    & \textbf{Lo que le preguntan} &
    \textbf{Lo que hay que resolver primero} \\
    \midrule
    \azul{Psicología} &
      ``¿Cómo es el perfil de los estudiantes que llegan a consulta?'' &
      ¿Perfil en qué variables: edad, semestre, motivo de consulta? ¿Y cómo se
      resume cada una? \\
    \addlinespace
    \azul{Psicología} &
      ``¿Cuánto tiempo espera un usuario para su primera cita?'' &
      ¿El tiempo \emph{típico} es el promedio, la mitad o el más frecuente?
      Los tres pueden dar cifras distintas \\
    \addlinespace
    \azul{T.\ Social} &
      ``¿El programa de acompañamiento familiar está funcionando?'' &
      ¿Funcionando en qué indicador, medido cómo y comparado con qué? \\
    \addlinespace
    \azul{T.\ Social} &
      ``¿A qué comuna dirigimos los recursos del próximo semestre?'' &
      ¿Qué datos tengo, sobre cuántos casos, de qué periodo y de qué fuente? \\
    \bottomrule
  \end{tabular}
  \end{center}

  \vspace{0.1cm}
  \begin{alertblock}{Lo que este curso le da}
    \tiny
    Ninguna de las cuatro se responde con una fórmula suelta. Se responden con
    un \destacar{procedimiento}: definir bien la pregunta, organizar los datos,
    resumirlos con la medida adecuada e interpretarlos \azul{sin decir más de
    lo que los datos permiten}.
  \end{alertblock}

  \vspace{0.05cm}
  {\tiny\color{uniGris} En psicología y trabajo social los datos casi nunca son
  dinero: son \destacar{personas, casos y trayectorias}. Eso no cambia el
  procedimiento, pero sí sube el costo de equivocarse: detrás de cada fila de
  la tabla hay alguien.}
\end{frame}

\begin{frame}{La prueba de las cuatro frases}
  \scriptsize
  \begin{center}
  \begin{tabular}{@{}p{7.8cm}p{4.2cm}@{}}
    \toprule
    \textbf{Afirmación} & \textbf{¿Se puede sostener?} \\
    \midrule
    ``Los 40 usuarios encuestados esperaron en promedio 12 días.'' &
      \azul{Sí}: habla de los medidos \\
    ``Los usuarios del centro esperan en promedio 12 días.'' &
      \destacar{No}: nadie midió a todos \\
    ``El 35\% de los encuestados asistió a todas las sesiones.'' &
      \azul{Sí} \\
    ``El 35\% de la población atendida asiste a todas las sesiones.'' &
      \destacar{No} \\
    \bottomrule
  \end{tabular}
  \end{center}

  \vspace{0.15cm}
  \begin{alertblock}{La señal que las distingue}
    \scriptsize
    Fíjese en el \destacar{sujeto de la frase}. Si dice ``los encuestados'',
    ``los casos registrados'' o ``la muestra'', la afirmación se sostiene con
    los datos que tiene. Si dice ``los usuarios'' a secas, \destacar{está
    hablando de gente que nunca midió}.
  \end{alertblock}

  \vspace{0.08cm}
  {\tiny\color{uniGris} La corrección casi siempre es de \azul{redacción}, no
  de cálculo: basta con nombrar bien a quién se midió. Es el error que más se
  penaliza en los tres exámenes.}
\end{frame}

\begin{frame}{El curso en tres unidades}
  \begin{columns}[T,onlytextwidth]
    \begin{column}{0.32\textwidth}
      \begin{block}{Unidad 1 (sem. 1--5)}
        \small Conceptos y tablas\\[0.2em]
        Población y muestra, tipos de variables, escalas de medición, tablas de
        frecuencia y métodos gráficos.
      \end{block}
    \end{column}
    \begin{column}{0.32\textwidth}
      \begin{block}{Unidad 2 (sem. 7--11)}
        \small Medidas estadísticas\\[0.2em]
        Tendencia central, posición, dispersión y forma; gráficos e
        interpretación de resultados.
      \end{block}
    \end{column}
    \begin{column}{0.32\textwidth}
      \begin{block}{Unidad 3 (sem. 13--15)}
        \small Análisis integrado\\[0.2em]
        Aplicación completa a un conjunto de datos y sustentación de decisiones
        con evidencia.
      \end{block}
    \end{column}
  \end{columns}

  \vspace{0.45cm}
  \begin{center}
    \small Curso \destacar{100\% presencial} · 3 créditos · 16 semanas ·
    2 horas de clase por semana
  \end{center}
\end{frame}

\begin{frame}{Evaluación: tres exámenes, nada más}
  \begin{center}
  \begin{tabular}{@{}llr@{}}
    \toprule
    \textbf{Evaluación} & \textbf{Semana y fechas} & \textbf{Peso} \\
    \midrule
    Examen 1 — Unidad 1 & Semana 6: 14 al 19 de septiembre & 35\% \\
    Examen 2 — Unidad 2 & Semana 12: 26 al 31 de octubre   & 35\% \\
    Examen 3 — Unidad 3 & Semana 16: 23 al 28 de noviembre & 30\% \\
    \midrule
    \textbf{Total} & & \textbf{100\%} \\
    \bottomrule
  \end{tabular}
  \end{center}

  \vspace{0.3cm}
  \begin{alertblock}{Cómo son los exámenes}
    \footnotesize
    Individuales, en \azul{Aulas Virtuales (Moodle)}, presentados de forma
    presencial en el aula y con tiempo límite. Selección múltiple con única
    respuesta y problemas de aplicación con respuesta numérica. Se permite
    calculadora y formulario.
  \end{alertblock}
\end{frame}

\begin{frame}{Acuerdos mínimos del curso}
  \begin{itemize}
    \item \azul{No hay talleres, quices ni trabajos calificables.} Todo lo que
          hagamos en clase es formativo: sirve para prepararse.
    \item Las semanas \destacar{6, 12 y 16} se dedican por completo al examen;
          no se dictan temas nuevos.
    \item Notas y retroalimentación se publican dentro de la semana siguiente
          al examen. Reclamaciones: 3 días hábiles.
    \item Presentación en fecha distinta solo con excusa válida radicada ante
          la coordinación del programa.
    \item Comunicación en clase, por correo institucional o por medio del
          representante de curso.
    \item Fraude o plagio anula la prueba y se reporta según el Reglamento
          Estudiantil.
  \end{itemize}
\end{frame}

\seccion{Unidad 1. Conceptos básicos y tablas de distribución}

\begin{frame}{1.1 Los cuatro conceptos de partida}
  \footnotesize
  \begin{center}
  \begin{tabular}{@{}p{2.4cm}p{5.0cm}p{4.0cm}@{}}
    \toprule
    \textbf{Concepto} & \textbf{Definición} & \textbf{Ejemplo} \\
    \midrule
    Población & Conjunto \emph{completo} de elementos sobre los que se quiere
      concluir. Tamaño $N$. & Los 1\,200 clientes de la tienda \\
    \addlinespace
    Muestra & Subconjunto de la población que sí se observa. Tamaño $n$. &
      40 clientes encuestados \\
    \addlinespace
    Variable & Característica que se mide y que \azul{puede cambiar} de un
      elemento a otro. & Gasto mensual, canal de compra \\
    \addlinespace
    Dato & El valor concreto que toma la variable en un elemento. &
      \$35\,000; ``App'' \\
    \bottomrule
  \end{tabular}
  \end{center}

  \vspace{0.2cm}
  {\footnotesize\color{uniGris} Si algo no cambia entre los elementos no es una
  variable, es una \azul{constante}: no aporta información para describir.}
\end{frame}

\begin{frame}{Población objetivo, población de estudio y marco muestral}
  \scriptsize
  \begin{center}
  \begin{tabular}{@{}p{3.0cm}p{4.2cm}p{4.6cm}@{}}
    \toprule
    \textbf{Concepto} & \textbf{Qué es} & \textbf{En el ejemplo} \\
    \midrule
    Población objetivo &
      Sobre quién \azul{quiero} concluir &
      Todos los clientes del almacén \\
    \addlinespace
    Población de estudio &
      A quién \azul{puedo} llegar de verdad &
      Los 1\,200 clientes \emph{registrados} \\
    \addlinespace
    Marco muestral &
      La \azul{lista} de la que se saca la muestra &
      La base de datos de registrados \\
    \addlinespace
    Muestra &
      Los que efectivamente se midieron &
      Los 40 encuestados \\
    \bottomrule
  \end{tabular}
  \end{center}

  \vspace{0.1cm}
  \begin{alertblock}{Dónde se cuela el sesgo}
    \scriptsize
    \destacar{En el salto de la primera fila a la segunda.} Los clientes que
    compran y nunca se registran \azul{no pueden salir en la muestra}, por más
    bien que se haga el muestreo. Si esos clientes son distintos —compran menos,
    o son de paso—, las conclusiones ya nacen torcidas.
  \end{alertblock}

  \vspace{0.05cm}
  {\tiny\color{uniGris} Por eso un informe serio \destacar{declara las tres
  cosas}: sobre quién quería concluir, a quién pudo llegar y a quién midió. Es
  la diferencia entre un estudio y una encuesta improvisada.}
\end{frame}

\begin{frame}{Unidad de observación y variable: no confundirlas}
  \begin{columns}[T,onlytextwidth]
    \begin{column}{0.48\textwidth}
      \begin{block}{Unidad de observación}
        \footnotesize
        El \azul{elemento} del que se toman los datos: un cliente, una factura,
        un empleado, un pedido.

        \vspace{0.1cm}
        Responde a ``¿de quién o de qué estoy midiendo?''.
      \end{block}
    \end{column}
    \begin{column}{0.48\textwidth}
      \begin{exampleblock}{Variable}
        \footnotesize
        La \azul{característica} que se mide en esa unidad: el gasto, la edad,
        el canal de compra.

        \vspace{0.1cm}
        Responde a ``¿qué estoy midiendo?''.
      \end{exampleblock}
    \end{column}
  \end{columns}

  \vspace{0.3cm}
  \begin{alertblock}{En una tabla de datos}
    \footnotesize
    Cada \destacar{fila} es una unidad de observación y cada
    \destacar{columna} es una variable. Es la estructura de un
    \azul{data.frame} en R, y la misma en cualquier software estadístico.
  \end{alertblock}
\end{frame}

\begin{frame}{De la realidad al dato: cómo se mide algo que no es un número}
  \scriptsize
  \begin{center}
  \begin{tikzpicture}[scale=0.9,every node/.style={font=\tiny,align=center},
    nodo/.style={draw=uniAzul,fill=uniAzul!10,rounded corners,inner sep=3pt,
                 minimum width=2.5cm,minimum height=0.62cm},
    flecha/.style={->,>=latex,uniGris,thick}]
    \node[nodo] (a) at (0,0) {Concepto\\ ``satisfacción''};
    \node[nodo] (b) at (3.4,0) {Variable\\ ``nivel de\\ satisfacción''};
    \node[nodo] (c) at (6.8,0) {Indicador\\ escala de 1 a 5};
    \node[nodo,fill=uniCrema,draw=uniTitulo] (d) at (10.2,0)
      {Dato\\ ``4''};
    \draw[flecha] (a)--(b); \draw[flecha] (b)--(c); \draw[flecha] (c)--(d);
  \end{tikzpicture}
  \end{center}

  \vspace{0.15cm}
  \begin{alertblock}{Las decisiones que se toman en el camino}
    \scriptsize
    ``Satisfacción'' no existe como número: \destacar{alguien decidió} medirla
    con una escala de 1 a 5, y no de 1 a 10, ni con una pregunta abierta.
    \azul{Esa decisión es parte del resultado.} Dos encuestas que preguntan
    distinto no son comparables, aunque las dos hablen de satisfacción.
  \end{alertblock}

  \vspace{0.1cm}
  \begin{exampleblock}{La pregunta que hay que hacerle a todo dato}
    \scriptsize
    \destacar{¿Cómo se produjo este número?} Un gasto de ``227 mil'' puede ser
    lo que el cliente \emph{dijo} que gasta, lo que la caja \emph{registró} o
    lo que el sistema \emph{estimó}. Los tres se ven igual en la tabla y
    significan cosas distintas.
  \end{exampleblock}
\end{frame}

\begin{frame}{Ejercicio 1 — en clase}
  \begin{alertblock}{Individual (6 minutos)}
    Un supermercado quiere estudiar el comportamiento de compra de sus
    \destacar{1\,200 clientes registrados}. Se encuesta a \destacar{40} de
    ellos y a cada uno se le pregunta el gasto del último mes y el canal por el
    que compró.

    \vspace{0.3cm}
    \begin{enumerate}
      \item Identifique la población y su tamaño.
      \item Identifique la muestra y su tamaño.
      \item ¿Cuál es la unidad de observación?
      \item Nombre las dos variables y dé un ejemplo de dato para cada una.
    \end{enumerate}
  \end{alertblock}
\end{frame}

\begin{frame}{Ejercicio 1 — solución}
  \begin{columns}[T,onlytextwidth]
    \begin{column}{0.50\textwidth}
      \begin{block}{Respuestas}
        \footnotesize
        \begin{itemize}
          \item Población: los 1\,200 clientes registrados, $N=1\,200$.
          \item Muestra: los 40 encuestados, $n=40$.
          \item Unidad de observación: \azul{un cliente}.
          \item Variables: \azul{gasto mensual} (dato: 35 mil) y \azul{canal de
                compra} (dato: ``App'').
        \end{itemize}
      \end{block}
    \end{column}
    \begin{column}{0.46\textwidth}
      \begin{exampleblock}{Fíjese}
        \footnotesize
        El número 1\,200 \azul{no} es un dato del estudio: es el tamaño de la
        población.

        \vspace{0.2cm}
        Y ``40'' tampoco: es el tamaño de la muestra. Los datos son los 40
        gastos y los 40 canales.
      \end{exampleblock}
    \end{column}
  \end{columns}
\end{frame}

\begin{frame}{Las etapas del proceso estadístico}
  \small
  \begin{enumerate}
    \item \azul{Plantear la pregunta.} ¿Qué se quiere saber y sobre quién?
    \item \azul{Definir población y unidad de observación.}
    \item \azul{Elegir las variables} y cómo se van a medir.
    \item \azul{Recolectar los datos} (censo o muestra).
    \item \azul{Organizar}: tablas de frecuencia. \hfill{\scriptsize (sesión 3)}
    \item \azul{Representar}: gráficos. \hfill{\scriptsize (sesión 4)}
    \item \azul{Resumir}: medidas estadísticas. \hfill{\scriptsize (unidad 2)}
    \item \azul{Interpretar y decidir.} \hfill{\scriptsize (unidad 3)}
  \end{enumerate}

  \vspace{0.2cm}
  \begin{alertblock}{Dónde se pierde la mayoría de los estudios}
    \footnotesize
    En los pasos 1 a 3. Si la variable está mal definida o la unidad de
    observación es ambigua, \destacar{ningún cálculo posterior lo arregla}.
  \end{alertblock}
\end{frame}

\begin{frame}{Estudio observacional o experimental}
  \scriptsize
  \begin{center}
  \begin{tabular}{@{}p{2.2cm}p{4.6cm}p{4.8cm}@{}}
    \toprule
    & \textbf{Observacional} & \textbf{Experimental} \\
    \midrule
    Qué se hace & Se \azul{mira} lo que pasa, sin intervenir &
      Se \azul{interviene} y se compara con un grupo de control \\
    \addlinespace
    Ejemplo & Registrar el gasto de los clientes que llegan &
      Poner una promoción en una sede y no en otra, y comparar \\
    \addlinespace
    Qué permite decir & Que dos cosas \azul{van juntas} &
      Que una \destacar{causa} la otra \\
    \addlinespace
    Costo y control & Barato, pero con muchas explicaciones alternativas &
      Costoso, a veces imposible o no ético \\
    \bottomrule
  \end{tabular}
  \end{center}

  \vspace{0.06cm}
  \begin{alertblock}{Por qué se lo digo en la primera semana}
    \tiny
    Porque el error más caro de todo el curso es leer un estudio
    \destacar{observacional} y redactar una conclusión \destacar{causal}.
    ``Los clientes que usan la app gastan más'' \azul{no} significa que la app
    los haga gastar más: puede que los que ya gastaban más sean justamente
    quienes se bajan la app.
  \end{alertblock}

  {\tiny\color{uniGris} Casi todo lo que veremos en este curso es
  observacional. Volveremos sobre esto en la unidad 3.}
\end{frame}

\begin{frame}{Ejercicio 2 — en clase}
  \begin{alertblock}{En parejas (8 minutos)}
    Para cada estudio, identifique la \azul{unidad de observación} y proponga
    \azul{dos variables} que tendría sentido medir.

    \vspace{0.25cm}
    \begin{enumerate}
      \item Un hospital quiere caracterizar las consultas del último mes.
      \item Una empresa quiere estudiar sus facturas de venta del año.
      \item Un colegio quiere describir a sus docentes.
      \item Una tienda quiere analizar sus pedidos a domicilio.
    \end{enumerate}
  \end{alertblock}
\end{frame}

\begin{frame}{Ejercicio 2 — solución}
  \scriptsize
  \begin{center}
  \begin{tabular}{@{}clll@{}}
    \toprule
    \textbf{\#} & \textbf{Estudio} & \textbf{Unidad de observación} &
    \textbf{Variables posibles} \\
    \midrule
    1 & Consultas del mes & Una consulta &
      Duración en minutos; especialidad \\
    2 & Facturas del año  & Una factura &
      Valor facturado; forma de pago \\
    3 & Docentes          & Un docente &
      Años de experiencia; nivel de formación \\
    4 & Pedidos a domicilio & Un pedido &
      Tiempo de entrega; canal del pedido \\
    \bottomrule
  \end{tabular}
  \end{center}

  \vspace{0.25cm}
  \begin{alertblock}{El error típico}
    \footnotesize
    Decir que la unidad del caso 1 es ``el hospital''. El hospital es
    \emph{uno solo}: si fuera la unidad, habría un solo dato y nada que
    describir. \destacar{La unidad es aquello de lo que hay muchos.}
  \end{alertblock}
\end{frame}

\begin{frame}{De dónde salen los datos}
  \scriptsize
  \begin{columns}[T,onlytextwidth]
    \begin{column}{0.48\textwidth}
      \begin{block}{Según quién los recogió}
        \tiny
        \textbf{Primarios:} los recoge usted para \azul{esta} pregunta.
        Encuesta, observación, experimento. Más costosos, pero a la medida.

        \vspace{0.1cm}
        \textbf{Secundarios:} ya existían. Registros de la empresa, DANE,
        bases públicas. Baratos, pero fueron pensados para otra pregunta.
      \end{block}
    \end{column}
    \begin{column}{0.48\textwidth}
      \begin{block}{Según cuánto se mide}
        \tiny
        \textbf{Censo:} se mide \azul{toda} la población. Exacto, pero
        costoso, lento y a veces imposible.

        \vspace{0.1cm}
        \textbf{Muestra:} se mide una parte. Más rápido y barato, y si está
        bien tomada, suficiente.
      \end{block}
    \end{column}
  \end{columns}

  \vspace{0.1cm}
  \begin{alertblock}{Por qué casi siempre se muestrea}
    \scriptsize
    Porque medir a los 1\,200 clientes cuesta 30 veces más que medir a 40, y
    porque a veces la medición \destacar{destruye} la unidad: nadie prueba la
    resistencia de los 10\,000 tornillos del lote, porque se quedaría sin
    lote.
  \end{alertblock}
\end{frame}

\begin{frame}{Cómo se arruina un estudio antes de calcular nada}
  \scriptsize
  \begin{center}
  \begin{tabular}{@{}p{3.2cm}p{4.2cm}p{4.6cm}@{}}
    \toprule
    \textbf{Problema} & \textbf{Ejemplo} & \textbf{Consecuencia} \\
    \midrule
    Muestra mal tomada &
      Encuestar solo los sábados en la mañana &
      Describe ese momento, no el negocio \\
    \addlinespace
    Pregunta ambigua &
      ``¿Cuánto gasta usted?'' (¿al día? ¿al mes?) &
      Datos que no son comparables entre sí \\
    \addlinespace
    Unidad mal definida &
      Contar ``clientes'' cuando el registro son facturas &
      Un cliente con 3 facturas cuenta tres veces \\
    \addlinespace
    No respuesta selectiva &
      Los insatisfechos no contestan &
      Todo se ve mejor de lo que es \\
    \bottomrule
  \end{tabular}
  \end{center}

  \vspace{0.1cm}
  \begin{alertblock}{La regla que hay que interiorizar}
    \scriptsize
    \destacar{Ningún cálculo arregla un dato mal recogido.} Se puede calcular
    la media de datos basura con seis decimales, y sigue siendo basura con seis
    decimales.
  \end{alertblock}
\end{frame}

\begin{frame}{Ejercicio 3 — en clase}
  \begin{alertblock}{Individual (8 minutos)}
    Una empresa con \destacar{850 empleados} encuesta a \destacar{120}. Parte
    de la tabla de datos se ve así:

    \vspace{0.2cm}
    \centering\footnotesize
    \begin{tabular}{@{}cccc@{}}
      \toprule
      \textbf{Empleado} & \textbf{Área} & \textbf{Antigüedad (años)} &
      \textbf{Salario (miles)} \\
      \midrule
      1 & Ventas       & 3 & 2\,100 \\
      2 & Operaciones  & 7 & 1\,850 \\
      3 & Ventas       & 1 & 1\,950 \\
      \bottomrule
    \end{tabular}

    \vspace{0.2cm}
    \raggedright
    \begin{enumerate}
      \item ¿Cuántas filas y cuántas columnas de variables tiene la tabla completa?
      \item Clasifique los números 850, 120, 7 y 1\,850: ¿cuál es dato, cuál es
            tamaño de población y cuál de muestra?
      \item ¿Es ``Empleado'' una variable?
    \end{enumerate}
  \end{alertblock}
\end{frame}

\begin{frame}{Ejercicio 3 — solución}
  \begin{columns}[T,onlytextwidth]
    \begin{column}{0.50\textwidth}
      \begin{block}{Puntos 1 y 2}
        \footnotesize
        \begin{itemize}
          \item \azul{120 filas} (una por empleado encuestado) y
                \azul{3 columnas} de variables.
          \item $N=850$ y $n=120$ son \azul{tamaños}, no datos.
          \item $7$ y $1\,850$ sí son \azul{datos}: valores de antigüedad y
                salario del empleado 2.
        \end{itemize}
      \end{block}
    \end{column}
    \begin{column}{0.46\textwidth}
      \begin{exampleblock}{Punto 3}
        \footnotesize
        No. ``Empleado'' es el \azul{identificador} de la unidad de
        observación: dice de quién es la fila, no qué se midió.
      \end{exampleblock}
      \begin{alertblock}{Ojo}
        \footnotesize
        Si el identificador fuera un número (1, 2, 3\dots), promediarlo daría
        un resultado sin ningún sentido. Es el mismo problema del código de
        sucursal que veremos en la sesión 2.
      \end{alertblock}
    \end{column}
  \end{columns}
\end{frame}

\begin{frame}{El conjunto de datos que usaremos todo el curso}
  \footnotesize
  Gasto mensual, en miles de pesos, de una muestra de \destacar{$n=40$}
  clientes del \azul{Almacén Bello}:

  \vspace{0.2cm}
  \begin{center}
  \ttfamily\scriptsize
  112\quad 127\quad 138\quad 149\quad 153\quad 158\quad 165\quad 171\quad 178\quad 184\\
  190\quad 195\quad 199\quad 203\quad 208\quad 212\quad 215\quad 219\quad 222\quad 226\\
  229\quad 231\quad 234\quad 238\quad 241\quad 245\quad 249\quad 253\quad 258\quad 264\\
  269\quad 275\quad 281\quad 288\quad 295\quad 302\quad 309\quad 318\quad 331\quad 352
  \end{center}

  \vspace{0.2cm}
  \begin{alertblock}{Guárdelo}
    \scriptsize
    Estos mismos 40 datos aparecerán en todas las sesiones: primero para armar
    tablas (sesión 3), luego gráficos (sesión 4), después medidas (unidad 2) y
    al final el análisis completo (unidad 3). \destacar{Así se ve cómo cada
    herramienta aporta algo distinto sobre el mismo conjunto.}
  \end{alertblock}
\end{frame}

\begin{frame}{La matriz de datos: cómo se guarda todo esto}
  \scriptsize
  \begin{center}
  \begin{tabular}{@{}crrll@{}}
    \toprule
    \textbf{Cliente} & \textbf{Gasto} & \textbf{Compras} &
    \textbf{Canal} & \textbf{Fidelidad} \\
    \midrule
    1 & 112 & 2 & Tienda & Nuevo \\
    2 & 127 & 3 & App    & Frecuente \\
    3 & 138 & 1 & Web    & Ocasional \\
    4 & 149 & 4 & Tienda & Frecuente \\
    \multicolumn{5}{c}{\dots} \\
    40 & 352 & 8 & App   & Frecuente \\
    \bottomrule
  \end{tabular}
  \end{center}

  \vspace{0.1cm}
  \begin{columns}[T,onlytextwidth]
    \begin{column}{0.48\textwidth}
      \begin{block}{Cómo se lee}
        \tiny
        Cada \azul{fila} es una unidad de observación (un cliente). \\
        Cada \azul{columna} es una variable. \\
        Cada \azul{celda} es un dato. \\
        $n$ = número de filas; aquí, 40.
      \end{block}
    \end{column}
    \begin{column}{0.48\textwidth}
      \begin{alertblock}{Lo que nunca debe pasar}
        \tiny
        Dos unidades en una misma fila; dos variables en una misma columna;
        una celda con \destacar{texto y número mezclados}
        (``\$127 mil aprox.''); o un dato faltante escrito como \azul{cero},
        que después se promedia como si fuera un gasto real.
      \end{alertblock}
    \end{column}
  \end{columns}

  \vspace{0.05cm}
  {\tiny\color{uniGris} Esta estructura es universal: es una hoja de cálculo,
  una tabla de base de datos y un \azul{data.frame} de R. \destacar{Dejar los
  datos así es la mitad del trabajo de cualquier análisis real.}}
\end{frame}

\begin{frame}{Ejercicio 4 — en clase (taller principal)}
  \begin{alertblock}{En grupos de 3 (20 minutos)}
    \scriptsize
    Su grupo asesora a la administración de la universidad, que quiere
    \destacar{reducir las filas en la cafetería}.

    \vspace{0.2cm}
    \begin{enumerate}\setlength{\itemsep}{-0.06cm}
      \item Formule la pregunta estadística en una sola frase.
      \item Defina la población y la unidad de observación. ¿Son estudiantes,
            visitas o días?
      \item Proponga \destacar{tres variables} y diga exactamente cómo se
            mediría cada una (con su unidad).
      \item ¿Censo o muestra? Justifique.
      \item Describa cómo tomaría los datos: quién, cuándo y dónde.
      \item Enuncie \destacar{un sesgo} que podría colarse en su diseño.
    \end{enumerate}

    \vspace{0.1cm}
    Al final, dos grupos exponen y el curso busca el sesgo del otro diseño.
  \end{alertblock}
\end{frame}

\begin{frame}{Ejercicio 4 — solución sugerida}
  \scriptsize
  \begin{columns}[T,onlytextwidth]
    \begin{column}{0.50\textwidth}
      \begin{block}{Puntos 1 a 3}
        \tiny
        \textbf{Pregunta:} ¿cuánto espera un usuario de la cafetería y de qué
        depende esa espera?

        \vspace{0.08cm}
        \textbf{Unidad:} \azul{una visita} a la cafetería (no un estudiante:
        un mismo estudiante puede ir varias veces).

        \vspace{0.08cm}
        \textbf{Variables:} tiempo de espera en minutos; franja horaria
        (mañana / mediodía / tarde); número de personas en fila al llegar.
      \end{block}
    \end{column}
    \begin{column}{0.47\textwidth}
      \begin{exampleblock}{Puntos 4 a 6}
        \tiny
        \textbf{Muestra:} un censo de todas las visitas de un semestre es
        inviable. Se mide una muestra de visitas.

        \vspace{0.08cm}
        \textbf{Recolección:} un observador con cronómetro, en \azul{días y
        franjas distintos}, durante dos semanas.

        \vspace{0.08cm}
        \textbf{Sesgo posible:} medir solo al mediodía —la hora pico— haría ver
        el problema \destacar{peor} de lo que es en promedio.
      \end{exampleblock}
    \end{column}
  \end{columns}

  \vspace{0.05cm}
  {\tiny\color{uniGris} No hay una única respuesta correcta. Lo que se evalúa
  es la \destacar{coherencia}: que la unidad, las variables y la forma de
  recolección respondan a la pregunta que se formuló en el punto 1.}
\end{frame}

\begin{frame}{Cómo se lee mal una estadística}
  \scriptsize
  \begin{center}
  \begin{tabular}{@{}p{5.0cm}p{6.8cm}@{}}
    \toprule
    \textbf{La afirmación} & \textbf{Lo que falta preguntar} \\
    \midrule
    ``El salario promedio subió 12\%'' &
      ¿Promedio o mediana? Unos pocos sueldos altos mueven el promedio \\
    \addlinespace
    ``El 80\% de los usuarios lo recomienda'' &
      ¿80\% \azul{de cuántos}? Ocho de diez no es lo mismo que 8\,000 de
      10\,000 \\
    \addlinespace
    ``Las ventas se dispararon'' (con gráfico) &
      ¿El eje vertical empieza en cero, o está cortado? \\
    \addlinespace
    ``Nueve de cada diez odontólogos lo recomiendan'' &
      ¿Quién los eligió? ¿Cuántos dijeron que no y no aparecen? \\
    \addlinespace
    ``Bajó la criminalidad un 3\%'' &
      ¿Bajaron los delitos o bajaron las \azul{denuncias}? \\
    \bottomrule
  \end{tabular}
  \end{center}

  \vspace{0.08cm}
  \begin{alertblock}{Las cuatro preguntas que desmontan casi cualquier cifra}
    \tiny
    \destacar{¿Sobre cuántos casos?} \quad \destacar{¿Cómo se midió?} \quad
    \destacar{¿Comparado con qué?} \quad \destacar{¿Qué no me están
    mostrando?}
  \end{alertblock}

  \vspace{0.05cm}
  {\tiny\color{uniGris} Este curso sirve para dos cosas: producir cifras
  honestas y \azul{no dejarse engañar por las ajenas}. La segunda le va a
  servir más veces en la vida que la primera.}
\end{frame}

\begin{frame}{Errores frecuentes en esta sesión}
  \scriptsize
  \begin{center}
  \begin{tabular}{@{}p{5.4cm}p{6.4cm}@{}}
    \toprule
    \textbf{Error} & \textbf{Corrección} \\
    \midrule
    Confundir la unidad de observación con la institución &
      La unidad es \azul{aquello de lo que hay muchos}: una consulta, una
      factura, un cliente \\
    \addlinespace
    Llamar ``dato'' al tamaño de la población &
      $N=1\,200$ describe el estudio, no es una observación \\
    \addlinespace
    Confundir la variable con el dato &
      La variable es ``gasto mensual''; el dato es ``227'' \\
    \addlinespace
    Decir ``los clientes'' cuando se midieron 40 &
      Se describe \azul{la muestra}: hay que nombrarla en la frase \\
    \addlinespace
    Empezar por el cálculo &
      Primero pregunta, población, unidad y variables \\
    \bottomrule
  \end{tabular}
  \end{center}
\end{frame}

\begin{frame}[fragile]{Primer vistazo a R: la tabla de datos}
  \scriptsize
  Todo lo de hoy tiene un nombre en R. La \azul{tabla de datos} es un
  \fx{data.frame}: una fila por unidad de observación, una columna por
  variable.
\begin{lstlisting}[style=Rmini]
clientes <- data.frame(
  id     = 1:4,
  gasto  = c(112, 127, 138, 149),      # variable cuantitativa
  canal  = c("Tienda", "App", "Web", "Tienda")   # variable cualitativa
)
clientes
nrow(clientes)   # cuantas unidades de observacion hay -> el tamano n
ncol(clientes)   # cuantas variables se midieron
\end{lstlisting}
\begin{lstlisting}[style=Rout]
  id gasto  canal
1  1   112 Tienda
2  2   127    App
3  3   138    Web
4  4   149 Tienda
[1] 4
[1] 3
\end{lstlisting}

  \vspace{-0.12cm}
  {\tiny\color{uniGris} \fx{nrow()} es el $n$ del estudio y \fx{ncol()} es el
  número de variables. En la sesión 14 volveremos sobre R con calma; por ahora
  basta ver que \destacar{la estructura es la misma} que dibujamos en el
  tablero.}
\end{frame}

\begin{frame}{Cierre de la sesión}
  \begin{block}{Lo que debe quedar claro hoy}
    \begin{itemize}
      \item Población es el todo ($N$); muestra es la parte observada ($n$).
      \item La unidad de observación es \emph{de quién} se mide; la variable es
            \emph{qué} se mide.
      \item Un dato es un valor concreto de una variable en una unidad.
      \item En una tabla: filas = unidades, columnas = variables.
    \end{itemize}
  \end{block}

  \vspace{0.25cm}
  \begin{alertblock}{Trabajo independiente para la próxima sesión}
    \footnotesize
    Resolver los \destacar{20 ejercicios de la lámina siguiente}; las
    respuestas están en la lámina posterior, para que usted mismo se corrija.
  \end{alertblock}
\end{frame}

\begin{frame}{Ejercicios propuestos}
  \vspace{-0.1cm}
  \begin{beamercolorbox}[wd=\textwidth,sep=0.10cm]{block body}
    \centering\scriptsize
    \textbf{1--10:} una universidad con 8\,000 estudiantes encuesta a 300 sobre
    su promedio acumulado y su programa académico.
  \end{beamercolorbox}
  \listatiny
  \vspace{0.15cm}
  \begin{columns}[T,onlytextwidth]
    \begin{column}{0.47\textwidth}
      \begin{enumerate}\setlength{\itemsep}{0.16cm}
        \item ¿Cuál es la población?
        \item ¿Cuál es $N$?
        \item ¿Cuál es la muestra?
        \item ¿Cuál es $n$?
        \item ¿Cuál es la unidad de observación?
        \item Nombre las dos variables
        \item Dé un dato posible de la primera variable
        \item Dé un dato posible de la segunda
        \item ¿Es 8\,000 un dato del estudio?
        \item ¿Cuántas filas tendría la tabla de datos?
      \end{enumerate}
    \end{column}
    \begin{column}{0.50\textwidth}
      \begin{enumerate}\setlength{\itemsep}{0.16cm}
        \setcounter{enumi}{10}
        \item ¿Cuántas columnas de variables tendría?
        \item ¿Población o muestra: los 40 clientes encuestados?
        \item ¿Variable o constante: el nombre del supermercado?
        \item ¿Qué describe la estadística descriptiva?
        \item ¿Sobre quiénes se puede afirmar algo con estos datos?
        \item ¿Se puede tener muestra sin población?
        \item Si se encuesta a todos, ¿cómo se llama el estudio?
        \item ¿Qué símbolo se usa para el tamaño poblacional?
        \item ¿Y para el tamaño muestral?
        \item En el conjunto del curso, ¿cuánto vale $n$?
      \end{enumerate}
    \end{column}
  \end{columns}
\end{frame}

\begin{frame}{Respuestas}{Autoevaluación de los ejercicios propuestos}
  \listatiny
  \vspace{0.15cm}
  \begin{columns}[T,onlytextwidth]
    \begin{column}{0.47\textwidth}
      \begin{enumerate}\setlength{\itemsep}{0.16cm}
        \item Los 8\,000 estudiantes
        \item $N=8\,000$
        \item Los 300 encuestados
        \item $n=300$
        \item Un estudiante
        \item Promedio acumulado y programa académico
        \item Por ejemplo, $4{,}1$
        \item Por ejemplo, ``Administración''
        \item No: es el tamaño de la población
        \item 300
      \end{enumerate}
    \end{column}
    \begin{column}{0.50\textwidth}
      \begin{enumerate}\setlength{\itemsep}{0.16cm}
        \setcounter{enumi}{10}
        \item 2
        \item Muestra
        \item Constante
        \item Lo que ya se observó
        \item Solo sobre los casos que se midieron
        \item No
        \item Censo
        \item $N$
        \item $n$
        \item $n=40$
      \end{enumerate}
    \end{column}
  \end{columns}
\end{frame}

\begin{frame}{Referencias}
  \footnotesize
  \begin{itemize}
    \item[{[1]}] Monroy Saldívar, S. (2008). \textit{Estadística descriptiva}.
      Instituto Politécnico Nacional.
      \texttt{https://elibro.net/es/lc/uniminuto/titulos/74722}
    \item[{[2]}] Ordóñez Fernández, F. F. \& González Fernández, J. (2021).
      \textit{Estadística descriptiva paso a paso}. Ediciones Pirámide.
      \texttt{https://elibro.net/es/lc/uniminuto/titulos/215449}
    \item[{[3]}] Llinás Solano, H. (2014). \textit{Introducción a la
      estadística con aplicaciones en ciencias sociales}. Universidad del
      Norte. \texttt{https://elibro.net/es/lc/uniminuto/titulos/70062}
  \end{itemize}

  \vspace{0.35cm}
  \begin{center}
    {\Large\color{uniAzul}\textbf{¿Preguntas?}}
  \end{center}
\end{frame}

\end{document}
